\exercise{Watching a kettle boil}{3}
Let's explore if observing the system can actually change its time evolution. 

\subquestion
Consider a diffusion process where a walker, moves randomly between the nodes of the linked-pair network so that jumps occur in average at rate 1. The walker starts at time zero in node 1. Compute the probability $x_1(2)$ that the walker is in node 1 at time 2.

\solution 
We already know that we can describe the diffusion on the linked-pair network as 
\eq{
\vec{\dot{x}} = -{\bf L} \vec{x}
}
where 
\eq{
{\bf L} = \avecc{1 & -1 \\ -1 & 1}
}
The eigenvectors of $\bf L$ are 
\eq{
\vec{v_1} = \avec{1\\ 1},\qqq \vec{v_2}=\avec{1\\ -1}.
}
whit the eigenvalues 
\eq{
\lambda_1=0, \qqq \lambda_2=2. 
}
Hence the general solution is 
\eq{
\vec{x}(t)=c_1\vec{v_1}+c_2\vec{v_2}\exp{-2t}
}
Now we consider a walker that is initially in the left node so 
\eq{
\vec{x}(0) = \avec{1 \\ 0 }
}
Which we can write as 
\eq{
\avec{1\\ 0} = \frac{1}{2} \avec{1\\ 1} + \frac{1}{2} \avec{1\\ -1} 
}
so that 
\eq{
c_1=\frac{1}{2}, \qqq c_2=\frac{1}{2}.
}
Hence we have the particular solution
\eq{
\vec{x}(t) = \frac{1}{2} \avec{1\\ 1} + \frac{1}{2} \avec{1\\ -1} \exp{-2t}
}
The question asks us to find $x_1(2)$ so we compute 
\eq{
R^* = x_1(2) = \frac{1}{2} + \frac{1}{2} \exp{-4} = \frac{1+\exp{-4}}{2} \approx 0.509
}
so slightly more than 50\%. We defined $R^*$ to be the result in this case to make it easier to talk about it later. 

\subquestion 
Now assume that we check on the walker at time 1. With probability $x_1(1)$ we find it in node 1, with probability $x_2(1)=1-x_1(1)$ we find it in node 2. Of course the measurement resets the probabilities. For example if we see the walker in node 1 at time 1, we can be sure that it is there and hence $x_1(1)$ becomes 1. 
Does knowing where the walker is at time 1 change the chance that it will be in node 1 at time 2? 

\solution 
\paragraph{Case 1.} Let's first consider the case where we check on the walker at time 1 and find it in node 1. The arrival of this new information puts us into the state 
\eq{
x_1(1) = \avec{1\\ 0}
}
which is just a way of saying that we are now certain that the walker is in node 1. Let's take a minute to appreciate that $\vec{x}$ does not describe the system, but rather our belief about the system. 

Let's plug the known state $x(1)$ into our general solution 
\eq{
\avec{1\\ 0} = c_1 \avec{1\\ 1} + c_2 \avec{1\\ -1}\exp{-2}
}
To accommodate this new situation we now need 
\eq{
c_1 = \frac{1}{2},\qqq c_2=\frac{1}{2}\exp{2}
}
So the arrival of the new information has changed our solution (!), it is now 
\eq{
\vec{x}(t) = \frac{1}{2}\avec{1\\ 1} + \frac{1}{2}\exp{2}\avec{1\\ -1} \exp{-2t}  
}
for $t\geq 1$. So we can now compute 
\eq{
R_1 = x_1(2) = \frac{1+\exp{-2}}{2} \approx 0.567
}
which is a bit more than before.

So in the case where we look at time 1 and see that the walker is in node 1, the chance of also finding the walker in node 1 at time 2 is greater than in the case where we did not look. 

\paragraph{Case 2.}So what about the case where we look at time 1 and we find the walker in node 2? 

Now we receive the new information 
\eq{
\vec{x}(1) = \avec{0\\ 1}  
}
and substituting this in the general solution we find 
\eq{
\avec{0\\ 1} = c_1 \avec{1\\ 1} + c_2 \avec{1\\ -1} \exp{-2}
}
so that we now need the coefficients 
\eq{
c_1 = \frac{1}{2},\qqq c_2 = - \exp{2} \frac{1}{2}
}
which give us the particular solution
\eq{
\vec{x}(t) = \frac{1}{2} \avec{1\\ 1}- \frac{1}{2} \exp{2} \avec{1\\ -1} \exp{-2t} 
}
for $t>1$. We can now compute 
\eq{
R_2=x_1(2) = \frac{1-\exp{2}}{2} \approx 0.432
}
So now there is slightly less than 50\% chance that the walker will be back in node 1 at time 2. 

\paragraph{The effect of observation.}So does observing the walker in itself change the result? Well, if we chose to observe the system at time 1, we don't know yet what result we will find so we have to average over both possibilities.
Using our particular solution from (a) we can compute the probability of finding the walker in node 1 when we look for it at time 1 
\eq{
p_1 = x_1(1) = \frac{1+\exp{-2}}{2} 
}
conversely the chance of finding the walker in node 2 when at time 1 is 
\eq{
p_2 = x_2(1) = \frac{1-\exp{-2}}{2} 
}
In the first case the probability of finding the walker in 1 at time 2 is 
\eq{
R_1 = \frac{1+\exp{-2}}{2}
}
In the second case the probability of finding the walker in 1 at time 2 is 
\eq{
R_2 = \frac{1-\exp{2}}{2} 
}
So in the scenario where we check on the system at time 1 but we do not know the result yet the expected result for the probainility of finding the walker in node 1 at time 2 is 
\eqa{
R&=& p_1 R_1 + p_2 R_2 \\
 &=& \frac{1+\exp{-2}}{2} \frac{1+\exp{-2}}{2} + \frac{1-\exp{-2}}{2} \frac{1-\exp{2}}{2}\\
 &=& \frac{1+2\exp{-2}+\exp{-4}}{4}+\frac{1-2\exp{-2}+\exp{-4}}{4} \\
  &=& \frac{2+2\exp{-4}}{4} \\
  &=& \frac{1+\exp{-4}}{2} = R^*
}
so we get exactly the same result as in the case where we did not check on the system at time 1. (I am so glad this worked out! Watching a kettle boil does not keep it from boiling.) 
